% Chapter 1 Introduction (untitled opening paragraphs)
% Funnel structure: broad context → specific problem → chapter roadmap → bridge to thesis

The nature of dark matter is an open question at the heart of fundamental physics.
Over the past century, independent observations have established that the visible matter we can detect --- stars, gas, dust --- accounts for only a small fraction of the total matter content of the Universe.
Flat rotation curves in spiral galaxies~\cite{Rubin:1970zza,Persic:1995ru}, the dynamics and gravitational lensing of galaxy clusters~\cite{Clowe:2006eq}, the detailed pattern of temperature anisotropies in the \CMB~\cite{Aghanim:2018eyx}, and the baryon acoustic oscillations imprinted on the large-scale distribution of galaxies~\cite{SDSS:2005xqv} all point consistently toward a dominant, non-luminous matter component.
These observations, spanning scales from individual galaxies to the observable Universe as a whole, converge on the $\Lambda$CDM model, in which approximately 85\% of the matter density is composed of cold, non-baryonic dark matter~\cite{Aghanim:2018eyx}.
The non-baryonic character of this component is independently confirmed by \BBN, which constrains the baryon density to $\Omega_b h^2 \approx 0.022$ --- far below the total matter density $\Omega_m h^2 \approx 0.14$ measured by the \CMB~\cite{Cirelli:2024ssz}.

If dark matter is not made of ordinary matter, then what is it?
Proposed candidates span nearly ninety orders of magnitude in mass, from ultralight axions to primordial black holes, and none is singled out by the data alone.
Among these, the \WIMP has long occupied a privileged position because of the so-called ``\WIMP miracle'': a stable particle with mass and coupling strength at the electroweak scale, produced thermally in the early Universe, naturally yields a relic abundance consistent with the observed dark matter density~\cite{Cirelli:2024ssz,Steigman:2012nb}.
The same annihilation process that sets the relic abundance also predicts that \glspl{WIMP_} should annihilate today in regions of high dark matter density, producing Standard Model particles --- gamma rays, neutrinos, positrons, and antiprotons --- that can in principle be detected by astrophysical observatories~\cite{Bergstrom:1997fj}.
This annihilation signature is what makes the \WIMP a natural target for the gamma-ray searches developed in this thesis, and it is why we adopt it as our working hypothesis.
% : not because it is the only viable candidate, but because its predicted Standard Model yields are what gamma-ray observatories can probe most directly.

This chapter lays the groundwork for the searches presented in the rest of the thesis.
We begin by reviewing the multi-scale evidence for dark matter, from galactic rotation curves to the \CMB (Section~\ref{sec:1.1}).
We then introduce the \WIMP paradigm, deriving the thermal freeze-out mechanism and the resulting relic abundance calculation that defines the target cross-section for experimental searches (Section~\ref{sec:1.2}).
Section~\ref{sec:1.3} surveys the three complementary experimental strategies --- direct, collider, and indirect detection --- and argues that indirect detection via gamma-rays provides the most direct test of the \WIMP hypothesis.
Finally, Section~\ref{sec:1.4} develops the indirect detection formalism in detail --- annihilation channels, spectral signatures, density profiles, and the $J$-factor/$D$-factor framework --- that connects particle physics parameters to observable gamma-ray fluxes.
The observational targets and the current status of indirect searches, reviewed at the end of this chapter, provide the direct motivation for the advanced statistical and machine learning approaches that constitute the core of this thesis.
